% !TeX root = ../nc-quadric-surface-avatar.tex
% 2026-11-22 / AS/avatar classification, source v7 Parts I-II
\section{AS正則\texorpdfstring{$\mathbb Z$}{Z}-代数とavatarの相互逆な分類}
\label{sec:qa-as-avatar-classification}

本節の目的は有限directed foundationの記述を
\emph{全整数頂点のAS正則$\mathbb Z$-代数の同型類}
との全単射へ昇格させることである。
元の証明は\cite[第2--8節，第11--12節]{AvatarVSeven2026}にある。
前二節で証明した曲線上の消滅・評価行列・canonical複体を入力とし、
残っていたAS--Gorenstein性、全helixの一意性と
必要条件の逆向きもここで示す。
一般的なspherical helix実現定理だけは
\cite[Theorem 5]{OkawaUedaSpherical2020}を外部入力とし、
本件でどのようにAS分解を得るかを明記する。

\subsection{最小分解で指定する二型}
\label{subsec:qa-as-minimal-types}

$A=\bigoplus_{j\ge i}A_{ji}$は局所有限、
$A_{ii}=k$、$A_{ji}=0$ ($j<i$)の
positively directed $\mathbb Z$-代数とする。
$P_i=e_iA$、$S_i$を右射影・右単純加群とする。
頂点をすべて保つ代数同型で分類し、
矢の基底や四周期同型の選択は追加データとしない。

\begin{definition}[二つのAS型]
\label{def:qa-as-square-rotated-minimal}
$I=\mathrm{rsq},\mathrm{sq}$に対し
$A\in\mathrm{AS}_{I}$とは、
すべての$i$で最小右射影分解
\begin{equation}
0\to P_{i-4}\to Q_{2,i}\to Q_{1,i}
 \to P_i\to S_i\to0
\label{eq:qa-as-two-minimal-resolutions}
\end{equation}
を持ち、さらに
\begin{equation}
\Ext_A^q(S_i,A)=0\ (q\ne3),\qquad
\Ext_A^3(S_i,A)\simeq {}_\ell S_{i-4}
\label{eq:qa-as-two-gorenstein}
\end{equation}
を満たすものとする。
$Q_{2,i},Q_{1,i}$は次の四周期表で指定する。
\[
\begin{array}{c|c|l|l}
I&i\bmod4&Q_{2,i}&Q_{1,i}\\ \hline
\mathrm{rsq}&0&P_{i-2}^{\oplus2}&P_{i-1}^{\oplus2}\\
&1&P_{i-3}^{\oplus2}&P_{i-2}^{\oplus2}\\
&2&P_{i-3}^{\oplus4}&P_{i-2}^{\oplus2}\oplus P_{i-1}^{\oplus2}\\
&3&P_{i-3}^{\oplus2}\oplus P_{i-2}^{\oplus2}&P_{i-1}^{\oplus4}\\ \hline
\mathrm{sq}&0&P_{i-3}^{\oplus2}\oplus P_{i-2}^{\oplus2}&P_{i-1}^{\oplus4}\\
&1&P_{i-2}^{\oplus2}&P_{i-1}^{\oplus2}\\
&2&P_{i-3}^{\oplus2}&P_{i-2}^{\oplus2}\\
&3&P_{i-3}^{\oplus4}&P_{i-2}^{\oplus2}\oplus P_{i-1}^{\oplus2}
\end{array}
\]
左右Noether性・幾何性・周期のtrivializationは仮定しない。
\end{definition}

\begin{lemma}[局所双対性と全$\mathbb Z$-代数の回収]
\label{lem:qa-as-serre-unique-recovery}
$A\in\mathrm{AS}_I$、$E_i=\pi P_i\in\Qgr A$とする。
$\mathcal D_{\rm lin}=\operatorname{thick}\{E_i\mid i\in\mathbb Z\}$
において連続する四対象はstrong exceptional tilting generatorである。
そのEnd代数$B$は大域次元二で、
Serre関手$S_B$により
\begin{equation}
E_{i+4}\simeq S_B^{-1}(E_i)[2],\qquad
A_{ji}=\Hom(E_i,E_j)\quad(j\ge i)
\label{eq:qa-as-serre-completion}
\end{equation}
を得る。全代数は初期四頂点のEnd代数から
頂点を保つ同型を除いて一意に決まる。
その数値Euler行列は
\[
G_{\rm rsq}=\begin{pmatrix}
1&0&2&6\\0&1&2&6\\0&0&1&4\\0&0&0&1
\end{pmatrix},\quad
G_{\rm sq}=\begin{pmatrix}
1&2&2&4\\0&1&0&2\\0&0&1&2\\0&0&0&1
\end{pmatrix},
\]
かつ$\dim A_{i+4m,i}=(2m+1)^2$である。
\end{lemma}
\begin{proof}
最小分解を各頂点で評価したEuler漸化式は、
負の成分が零で対角が$k$という初期条件から
全Hom次元を一意に定める。
可換基準collection
$(\OO,\OO(1,0),\OO(0,1),\OO(1,1))$と
$(\OO,\OO(-1,1),\OO(0,1),\OO(1,2))$
にも同じ漸化式が成立し、上の行列を得る。
四周期のHilbert関数は
$h^0(\PP^1\times\PP^1,\OO(2m,2m))=(2m+1)^2$。
これを可換基準モデルとの類比だけに依存させず、
最小分解表から有理母関数として直接検算する。
$r=i\bmod4$、$h_{r,n}=\dim A_{i,i-n}$、
$H_r(z)=\sum_{n\ge0}h_{r,n}z^n$とおき、
$Q_{1,i},Q_{2,i}$中の$P_{i-k}$の重複度を
それぞれ$a_{r,k},b_{r,k}$とする。
表のEuler漸化式は
\[
 (1-z^4)H_r(z)
 =1+\sum_{k=1}^3
 (a_{r,k}-b_{r,k})z^kH_{r-k\bmod4}(z).
 \tag{Euler4}\label{eq:qa-as-four-phase-euler}
\]
これら四本の一次方程式の有理関数解について、
$\zeta^4=1$を満たす原始四乗根$\zeta$を取ると、
各$r$で
\begin{equation}
 \frac14\sum_{j=0}^3H_r(\zeta^jz)
 =\frac{1+6z^4+z^8}{(1-z^4)^3}
 =\sum_{m\ge0}(2m+1)^2z^{4m}.
 \label{eq:qa-as-four-phase-rational-proof}
\end{equation}
左辺は四の倍数次数への射影である。
式の有理関数恒等性は

\noindent\texttt{computations/verify\_as\_avatar\_euler.py}で
四元連立の厳密な逆行列計算として再現できる。
ここで確認したのは\emph{最小分解が完全である場合の
Euler数値帰結}であって、
分解の完全性やAS--Gorenstein性を
母関数だけで証明したのではない。

まず有限区間で局所双対の消滅を確認する。
$P_j$ を頂点 $b$ 未満の部分加群で割った有限長商を
$U_{j,b}$ とする。頂点によるフィルトレーションの
因子は $S_\ell^{\oplus\dim A_{j\ell}}$（$b\le\ell\le j$）である。
\eqref{eq:qa-as-two-gorenstein} を射影 $P_i$ の成分で読むと、
$\Ext_A^q(S_\ell,P_i)$ は $(q,\ell)=(3,i+4)$ の場合だけ
$k$ となる。従ってフィルトレーションの長完全列から
$\Ext_A^q(U_{j,b},P_i)$ は $q\ne3$ では零であり、
$q=3$ の次元は $b\le i+4\le j$ なら $\dim A_{j,i+4}$、
それ以外は零である。
さらに切断を下げるとき、追加される因子の
$\Ext_A^3(-,P_i)$ は零なので、$b\le i+4$ 以後の
遷移写像は同型となる。
従ってこの計算はフィルトレーションの Euler 特性だけでなく、
導来 torsion を計算する直極限の各頂点での消滅を与える。

加群としての同定には、有限長加群 $N$ を
頂点 $i+4$ 未満の部分と $i+4$ 以下の部分で切る。
その二つの外側の部分の $\Ext_A^\bullet(-,P_i)$ は上の計算で零である。
従って包含と商の誘導する同型により
$\Ext_A^3(N,P_i)\simeq\Hom_k(N_{i+4},k)$ を得る。
この同型は $\Ext_A^3(S_{i+4},P_i)\simeq k$ の同定を固定すれば
$N$ に自然である。
特に $U_{j,b}$ に適用した同型は $j$ を変える道の作用と可換であり、
直極限で $R\tau(P_i)\simeq D(Ae_{i+4})[-3]$ となる。
この計算を tails に移すため、$J_{i,b}$ を
$P_i$ の頂点 $b$ 未満の部分加群とする。
$P_i/J_{i,b}$ は有限長で、各単純の長さ三の解消から
有限生成射影による長さ高々三の解消を持つ。
$J_{i,b}$ も有限生成射影による長さ高々二の解消を持つ。
さらに上の添字規約では、$J_{i,b}$ の射影解消を
頂点 $b$ 未満の射影の直和で選べる。
従って下に有界な torsion 加群 $N$ に対して、
十分小さい $b$ で $\Ext_A^p(J_{i,b},N)=0$ が全 $p$ で成り立つ。
一般の torsion 加群はそのような部分加群の filtered colimit
であり、有限射影解消により同じ極限消滅が成り立つ。
従って Gabriel 商の section 関手には
論文Bの \cite[式 \texttt{eq:cubic-tails-section-colimit} と
\texttt{eq:cubic-tails-derived-localization}]{QuadricModuli2026}
と同じ局所化三角を適用できる。
ここで用いたのは有限区間の射影解消と torsion 上の消滅である。

この三角を $P_j$ に適用して頂点 $i$ で評価すると、
\[
 \Hom(E_i,E_j)=A_{ji},\qquad
 \Ext^2(E_i,E_j)=D(A_{i,j+4}),\qquad
 \Ext^q(E_i,E_j)=0\quad(q>0,\ q\ne2)
\]
を得る。特に長さ四の窓の全対で高次 Ext が消える。
最小分解を$\pi$で移した四項完全列を
正負双方へ繰り返すと連続四対象は全$E_j$を生成する。
有限quiverの最小関係はsourceからsinkへ向かい、
関係同士のoverlapを持たない。
よってその有限End代数の単純分解の長さは高々二。
$E_j$ を末項とする四項窓の直和を $T$ とする。
$B=\End(T)$ は有限次元かつ有限大域次元なので、
導来 Morita 同値で Serre 関手は $DB$ との導来テンソルに対応する。
この窓の各 $E_i$ について $\Hom(E_i,E_{j-4})=0$ であり、
先の局所双対から
\[
 \RHom(E_i,E_{j-4}[2])\simeq D\Hom(E_j,E_i)[0]
\]
となる。この式は $i$ の窓内の射に自然である。
実際、$\Ext^2$ の同定は有限切断の $i$ 番射影への
道の作用を双対化したものである。
従って $\RHom(T,E_{j-4}[2])$ は、ただ一つのコホモロジーとして
右 $B$ 加群 $D\Hom(E_j,T)$ を持つ。
これは $S_B(E_j)$ に対応する加群であり、
Morita 同値により $S_B(E_j)\simeq E_{j-4}[2]$ を得る。
窓を変えても生成する圏は同じなので、その Serre 関手は同一である。
よって全 $j$ で $S_B^{-1}[2](E_j)\simeq E_{j+4}$ となる。
有限窓の導来圏でこの自己同値を反復し、合成を通常の射の合成に
取れば、全整数の代数の復元 \eqref{eq:qa-as-serre-completion} を得る。
これは周期同型の選択を分類データへ加える主張ではない。
\end{proof}

\subsection{AS条件からavatarを復元する}
\label{subsec:qa-as-to-avatar}

\begin{lemma}[AS条件が強制する階数]
\label{lem:qa-as-required-ranks}
rotated型の有限quiverを
$0\overset{2}{\longrightarrow}2
 \overset{4}{\longrightarrow}3$、
$1\overset{2}{\longrightarrow}2$と書き、
$R_0,R_1$を二つの二次元関係空間とする。
$V=H^0(2\to3)$とし、
二つの関係評価行列を$M_0(v),M_1(v)$、
随伴関係の$4\times4$行列を
$T(c,d),\widetilde T(a,b)$とすると
\begin{enumerate}
\item $R_0,R_1\subset V\otimes k^2$に
 非零の分解可能テンソルはない。
\item $\det M_0(v)=\det M_1(v)=0$、
 $v\ne0$なら両行列の階数は一。
\item $c,d$がともに非零なら
 $\operatorname{rk}T(c,d)\ge3$、
 $a,b$がともに非零なら
 $\operatorname{rk}\widetilde T(a,b)\ge3$。
\end{enumerate}
\end{lemma}
\begin{proof}
まず無限の関係を有限窓から回収する。
有限窓の代数 $B$ は大域次元二であり、$R_0,R_1$ は
有向道代数の関係 ideal の最小生成である。
未変形の Keller の relation completion
\cite[Theorem 6.10]{KellerDeformedCY2011} を用いると、
$T_B\RHom_{B^e}(B,B^e)[2]$ の零次コホモロジーは
有限窓の関係に双対戻り矢を加えた通常の Jacobian 代数である。
ここでは元の矢のテンソル次数を零、戻り矢を一とする。
この次数は一窓をまたぐ回数であり、関係の巡回微分は
両辺で同じ回数を保つ。
逆 Serre 双加群 $\RHom_{B^e}(B,B^e)[2]$ は
$S_B^{-1}[2]$ を表すので、その $m$ 次テンソル成分の
頂点対 $(i,j)$ は前補題により $\Hom(E_i,E_{j+4m})$ になる。
積は双加群のテンソル積、すなわち添字を移した射の合成である。
従ってこの Jacobian 代数を四頂点ずつ展開すると元の $A$ を得る。
以下の周期商で使う関係はこの通常の道 ideal の全生成である。

四周期の有限quiverのsourceとsinkを連結し、
各ブロック内の元の関係とその随伴微分関係を
すべて課した商を考える。
一ブロックのboundary cornerを$W$とすると、
$m$ブロックのcornerは$W^{\otimes m}$である。
境界間の道は一意に分割される。
以下では各場合ごとに境界を選び、
境界をまたぐ関係の像が零であることを確認する。
そのとき残る関係は各ブロックの一因子だけに作用するので、
ブロックごとの商のテンソル積が全境界 corner となる。
$\dim W\ge2$なら成分次元は少なくとも$2^m$となり、
前補題の$(2m+1)^2$に反する。

使う関係を先に書く。
$v_0,\dots,v_3$を$V$の基底とし、
$r_{0,\alpha}=\sum_{\ell,i}A_{\ell i\alpha}v_\ell a_i$、
$r_{1,\beta}=\sum_{\ell,j}B_{\ell j\beta}v_\ell b_j$
（$\alpha,\beta=0,1$）とする。
次周期の戻り矢$c_\alpha,d_\beta$を加えると、
各$r_{0,\alpha},r_{1,\beta}$に加えて
\[
\begin{aligned}
 \sum_{\ell,\alpha}A_{\ell i\alpha}c_\alpha v_\ell&=0
 &&(i=0,1),\\
 \sum_{\ell,\beta}B_{\ell j\beta}d_\beta v_\ell&=0
 &&(j=0,1),\\
 \sum_{i,\alpha}A_{\ell i\alpha}a_i'c_\alpha+
 \sum_{j,\beta}B_{\ell j\beta}b_j'd_\beta&=0
 &&(\ell=0,1,2,3)
\end{aligned}
\tag{Adj-rel}\label{eq:qa-as-forcing-adjoint-relations}
\]
および周期移動を持つ。以下は消した矢で切った
\emph{全随伴関係の像}を含む商である。

(1) 分解可能な関係を$r_{0,0}=v_0a_0$へ正規化する。
他方の枝の頂点 $1$ とその周期移動を零にし、
$v_0=a_0=c_1=0$を課す。境界は各周期の頂点 $3$ とする。
式~\eqref{eq:qa-as-forcing-adjoint-relations}の
$c_0$係数は$v_0$か$a_0$を含み、
$c_1$係数は零となるため、境界をまたぐ $c v$ 型の関係は消える。
残る境界から次の境界への道は三本で、
$r_{0,1}'c_0$による関係は高々一本。
従って$\dim W\ge3-1=2$である。
$R_1$の場合は二枝を交換する。

(2) $M_0(v)=0$かつ$\det M_1(v)=0$とする。
$v=v_0$とし、$M_1(v_0)$の左右kernelに
沿って$b_0,d_0$を選ぶ。
$v_1=v_2=v_3=b_1=d_1=0$を課し、境界を各周期の頂点 $2$ とする。
境界をまたぐ $v a,v b$ 型の $r_0,r_1$ はこの切断で消える。
$c v,d v$ 型の関係も、$M_0(v_0)=0$ と
$d_0$ の kernel 条件により消える。
$a'c,b'd$ 型のうち $\ell=0$ の式は $b_0$ の kernel 条件から零となる。
四本の$a_i'c_\alpha$と一本の$b_0'd_0$、
合計五本の境界道に、
$\ell=1,2,3$の高々三関係だけが残る。
よって$\dim W\ge5-3=2$である。
$M_1(v)=0$の場合も同様。

(3) $\operatorname{rk}T(c,d)\le2$なら
$c=d=(1,0)$に正規化し、
$A_{\ell i0},B_{\ell j0}$の四係数列が
$\langle v_0,v_1\rangle$に入るようにできる。
境界を各周期の頂点 $3$ とし、$v_0=v_1=c_1=d_1=0$を課す。
境界をまたぐ $c v,d v$ 型の関係は係数列が
$\langle v_0,v_1\rangle$ に入るため消える。
残った$a_0'c_0,a_1'c_0,b_0'd_0,b_1'd_0$を
高々二関係で割る空間$U'$は
$\dim U'\ge2$である。
次の境界までの道空間は
\[
 W=(\langle v_2',v_3'\rangle\otimes U')/
       \langle r_{0,1}'c_0,r_{1,1}'d_0\rangle
\]
であるから$\dim W\ge2\dim U'-2\ge2$。
$c_0,d_0$を可逆元や恒等射と扱っていないので、
この次元計算は四周期の正当な次数付き商に対するもの。
各場合で$\dim W\ge2$は指数増大に矛盾する。
左最小分解に同じ議論を適用すると
$\widetilde T$の階数条件を得る。
\end{proof}

\begin{proposition}[一点表示の幾何]
\label{prop:qa-as-extraction}
任意の$A\in\mathrm{AS}_{\rm rsq}$から
\[
Y=\PP_{\PP^1}(\mathcal E)\simeq\Sigma_0\ \text{または}\ \Sigma_2,
\quad
D\hookrightarrow Y,\quad
P\in\Pic^0(D)
\]
を回収できる。曲線型では
$D\in|2h|$、$\Sigma_2$なら$D\cap S=\varnothing$、
かつ$P\ne\OO_D,\OO_D(h-2f)$となる。
初期四項は
\[
(P,\OO_D(h-2f),\OO_D(h-f),\OO_D(2h-f)).
\]
$D=Y$の場合には$\Sigma_0$の古典型一点を得る。
\end{proposition}
\begin{proof}
保持三頂点の安定表現は
$\{([b],[v])\in\PP^1\times\PP^3\mid M_1(v)b=0\}$
である。$[b]$ を固定すると、関係は $V$ の二つの一次式となる。
この二式が従属なら、$R_1$ の非零元の $b$ による縮約が零となり、
その元は $V\otimes\ker b$ に属する。
$\ker b$ は一次元なので、これは分解可能テンソルを与え、
前補題に反する。
よって普遍関係写像
$\OO(-1)^2\to V\otimes\OO_{\PP^1}$
は各点で階数二であり、単元 minor により局所分裂単射である。
その余核 $\mathcal E$ は rank 二、次数二の大域生成束であり、
$M_1(v)b=0$ の incidence は $\PP(\mathcal E)$ と
試験スキーム上の商線束の関手として同型である。
$\mathcal E$ の分裂次数は大域生成性から非負、和は二なので
$\OO(1)^2$ または $\OO\oplus\OO(2)$ となり、
$Y$ の二型を得る。
忘れた頂点を加える評価は
$\kappa:\OO_Y^2\to\OO_Y(h)^2$である。
零点では $\det M_0(v)=\det M_1(v)=0$ であり、
前補題により $M_0(v)$ の階数は一なので、
局所的に$\kappa\sim\operatorname{diag}(1,q)$。
従って$\PP_Y(\coker\kappa)=V(q)=D$は
scheme-theoretic閉埋め込みであり、
$\det\kappa$の類は$2h$。
$\Sigma_2$ の負断面 $S$ は
$\mathcal E=\OO\oplus\OO(2)\twoheadrightarrow\OO$ の商に対応する。
この商では $\OO_Y(h)|_S\simeq\OO_S$ であり、
普遍評価 $V\otimes\OO_S\to\OO_S$ は一つの固定商 $[v_*]$ となる。
従って incidence 式 $M_1(v_*)b=0$ はすべての $[b]\in\PP^1$ で成り立つ。
二つの独立な $b$ を代入すると $M_1(v_*)=0$ を得る。
もし $\det M_0(v_*)=0$ なら、両行列式が零なのに
$M_1$ が階数零となり、前補題の第二条件に反する。
従って $\det M_0$ は $S$ に沿って単元であり、
スキームとして $D\cap S=\varnothing$ である。
この議論は $\Sigma_2$ 上で $\det\kappa$ が恒等的に零となる場合も排除する。
普遍右表現の双対と余核$Q$を用いて
$P=Q^{-1}\OO_D(2h-f)$と正規化する。

点表示の移動を随伴関係
\eqref{eq:qa-as-forcing-adjoint-relations} で記述する。
既知の評価を $v,a,b$ とすると、$c,d$ は
$M_0(v),M_1(v)$ の関係空間側の核線から得る。
いずれも階数一なので核線は局所直和因子であり、
単元 minor による記述は基底変更と可換である。
$c,d$ を代入した $T(c,d)$ では、既知の $v$ が左核に入るため
行列式は零となり、前補題の第三条件によって階数は三である。
従って右核から次の組 $(a',b')$ が一意に定まる。
その一方が零なら、他方と $c$ または $d$ によって
非零の関係テンソルの縮約が零となる。
この関係は二次元矢空間の一次元部分に支持されるので
分解可能となり、前補題の第一条件に反する。
従って $a',b'$ はともに非零である。

次に $\widetilde T(a',b')$ には既知の $(c,d)$ が核を与える。
同じ階数三の議論により、新しい $v'$ が一意に定まる。
随伴関係は $M_0(v')a'=M_1(v')b'=0$ を含むので、
$(v',a',b')$ は再び元の四頂点の点表示を与える。
逆向きにも同じ核線を用いると、階数一と階数三での
一意性により元の表示に戻る。
すべての写像は単元 minor 上の核線で構成されるので
逆性はスキーム上で成立する。
この移動の自己同型を $\sigma$ とし、
矢を $L_i\to L_j$ の評価と書くために普遍右表現を双対化する。
初期の線束を $L_0=P,L_1=HF^{-2},L_2=HF^{-1},L_3=H^2F^{-1}$
と正規化すると、二つの行列の余核は $Q$ および $HF$ である。
後者は $b:B\otimes\OO\twoheadrightarrow F$ の核に $H$ を捩った
部分束が $M_1(v)$ の像になることから従う。
したがって二つの階数一行列の核は、
行列式を取るとそれぞれ $QH^{-2}$ と $FH^{-1}$ である。
その双対が $L_4L_3^{-1}$ と $L_5L_3^{-1}$ を与えるので、
$L_4=H^2P$、$L_5=H^3F^{-2}$ を得る。

続く二つの階数三の写像では、
核と余核の線束を含む完全列の型は
\[
\begin{aligned}
0&\to L_6^{-1}\to
 (L_4^{-1})^{\oplus2}\oplus (L_5^{-1})^{\oplus2}
 \xrightarrow{T(c,d)} (L_3^{-1})^{\oplus4}
 \to L_2^{-1}\to0,\\
0&\to L_7^{-1}\to (L_6^{-1})^{\oplus4}
 \to (L_4^{-1})^{\oplus2}\oplus (L_5^{-1})^{\oplus2}
 \to L_3^{-1}\to0 .
\end{aligned}
\]
最初の余核は既知の $v$、次の余核は既知の $(c,d)$ の評価である。
二次元・四次元の定数係数空間の行列式を除けば、
$L_6=L_4^2L_5^2L_2L_3^{-4}=H^3F^{-1}P^2$ と
$L_7=L_6^4L_3L_4^{-2}L_5^{-2}=H^4F^{-1}P^6$ となる。
よって移動後の偏極は
$\sigma^*F=L_6L_5^{-1}=FP^2$ と
$\sigma^*H=L_7L_6^{-1}=HP^4$ である。
曲線型では、二つの式から $\sigma^*(HF^{-2})=HF^{-2}$ となる。
$\Sigma_0$ 上で成分の類を $af+bg$ と書くと、
$HF^{-2}$ のその成分上の次数は $a-b$ である。
この値が同じで類が異なる二成分が共存すると仮定すると、
一方の二係数は他方よりともに少なくとも一大きい。
小さい方も非零の有効類なので、二成分の和のいずれかの
係数が三以上となり、$D\in|2f+2g|$ に反する。
従って $\sigma$ は各成分の数値類を保ち、
特に $F$ の多重次数を保つ。
$\Sigma_2$ では各成分が負断面を避けるため、その類は
正整数倍の $h$ であり、類の異なる二成分は $2h$ 内に共存しない。
この場合も $F$ の多重次数は保たれる。
従って $\sigma^*F=FP^2$ を各成分で比較して $\deg P=0$ を得る。
ここで次数は被約既約成分への制限で測る。
多重次数零が $\Pic^0(D)$ を与える理由は、
まず $D_{\mathrm{red}}$ の正規化に移すと、
各正規化成分の Picard 群の単位成分が次数零部分だからである。
正規化による Picard 群の核は、導手の有限スキーム上の
単元群の商として記述され、連結である。
実際、代数閉体上の有限局所代数の単元群は
剰余体の乗法群と冪零イデアルに由来する連結な群からなる。
さらに $D$ の冪零イデアル $\mathcal J$ に対し
$1+\mathcal J$ を冪で濾過すると、逐次商は加法的な準連接層
$\mathcal J^r/\mathcal J^{r+1}$ である。
曲線上でこれらの $H^2$ は零で、$H^1$ はベクトル群なので、
被約化への Picard 制限は全射で核は連結である。
従って $P$ は非被約の場合も $\Pic^0(D)$ に属する。
曲線型では $\kappa$ の行列式は非零で、
曲面上の二項解消の余核は $\iota_*Q$ である。
これを $-2h+f$ と $-h-f$ で捩ると、余核はそれぞれ
$\iota_*P^{-1}$ と $\iota_*(HF^{-2}P^{-1})$ となる。
$\Sigma_0$ では $h=f+g$ なので、一つ目の捩れの
二つの曲面項はそれぞれ $-f-2g,-g$、
二つ目は $-2f-g,-f$ であり、Künneth によってすべて acyclic である。
$\Sigma_2$ では、相対次数 $-1$ の項の順像は全次数で零であり、
$-f$ の順像は $\OO_{\PP^1}(-1)$ である。
残る $-2h+f$ については相対双対性と
$\omega_\pi=\OO_Y(-2h+2f)$ により
$R^1\pi_*\OO_Y(-2h+f)=\OO_{\PP^1}(-1)$、$\pi_*=0$ となる。
従って両曲面で
$H^\bullet(P^{-1})=H^\bullet(HF^{-2}P^{-1})=0$ を得る。
自明線束には単位切断があるので、二つの除外条件が従う。
$\det\kappa\equiv0$ なら、負断面での可逆性により
$\Sigma_2$ は起こらない。
$Y=\Sigma_0$ 上で $\kappa$ は至るところ階数一なので、
余核 $Q$ は $Y$ 上の線束であり、点表示のスキームは $Y$ 全体である。
この場合の四周期の自己同型 $\sigma$ は
$\omega_Y^{-1}\simeq H^2$ を保つ。
$\Pic(Y)\simeq\mathbb Z^2$ には二捩れがないので、
$H$ は反標準束の一意な平方根として $\sigma$ に保たれる。
従って四周期の式 $\sigma^*H=HP^4$ から $P^4=\OO_Y$ を得て、
同じ Picard 群の計算から $P=\OO_Y$ となる。
この場合の線形代数には曲線型の単射解消を使わない。
\end{proof}


\subsection{Avatarからの全AS代数の逆構成}
\label{subsec:qa-as-avatar-inverse}

\begin{lemma}[canonical truncation]
\label{lem:qa-as-canonical-truncation}
rotated square avatarから前節の$L_i$を全整数に作る。
任意の$i$でcanonical四項列を
初期四頂点の共変Yonedaへ制限すると、
頂点単純加群の最小射影分解になる。
特に三回の右mutationは
\begin{equation}
T^{-1}_{L_{i+3}}T^{-1}_{L_{i+2}}T^{-1}_{L_{i+1}}(L_i)
 \simeq L_{i+4}[-2]
\label{eq:qa-as-canonical-mutation}
\end{equation}
を満たす。
\end{lemma}
\begin{proof}
$pencil$の評価列で作った
\eqref{eq:qa-core-rsq-complex0}と
\eqref{eq:qa-core-rsq-complex1}は
Yonedaで最初の二種の最小関係と一次syzygyを与える。
後二列は四行四列の対角分解の
Cartier曲線への導来制限と双対化から来るので
Tor項を含めて完全である。
各位相で最初に現れる非零な矢空間への評価は、
pencil または四次線束の全切断空間を使うので全射である。
第一・第二位相の残る頂点では、第一 pencil の評価列を
その頂点の線束で双対化して用いる。
評価の余核は pencil trick により次数二の線束の $H^1$ に
含まれるので零となる。第三位相の後二頂点では、四次線束 $H$
に対して $a_+$ または $b_+$ の評価列を用いる。
乗法の余核は $H a_+^{-1}$ または $H b_+^{-1}$ の $H^1$ に
含まれ、これらも次数二の pencil なので零である。
第四位相の最後の頂点では、第四列を作る前の曲面上の
$0\to\OO^{\oplus4}\to\OO(a)^{\oplus2}\oplus\OO(b)^{\oplus2}
 \to\iota_*Q\to0$
の大域切断が、$H^1(\OO)=0$ により所要の全射を与える。
中間の完全性もEuler特性だけには依存しない。
canonical複体を一時的に
$0\to L_i\to M_1\to M_2\to L_{i+4}\to0$と書く。
共変Yonedaの頂点$b$における
$\Hom(M_1,L_b)\to\Hom(L_i,L_b)$のkernelの次元を
計算すると
\[
\begin{array}{c|c|ccc}
i\bmod4&\text{相対次数}&b=1&b=2&b=3\\ \hline
0&(0,0,2,6)&0&0&2\\
1&(0,2,6,8)&0&2&4\\
2&(0,4,6,6)&0&2&2\\
3&(0,2,2,4)&0&0&4
\end{array}
\]
となる。ここで $L_b$ は窓の第 $b$ 頂点の略記である。
窓内の逆向きの通常の Hom は、
命題~\ref{prop:qa-core-forward-ext-vanishing}と
$\omega_D\simeq\OO_D$ による Serre 双対で零である。
従って通常の Hom を使ったこの計算は directed Yoneda の評価と一致する。
同じ双対性により、窓の外の端についても
$\Hom_D(L_{i+4},L_b)=\Ext_D^1(L_b,L_{i+4})^\vee=0$ となる。
特に総次数が等しい対でも、前向き消滅の次数零の場合がこれを保証する。

短完全列 $0\to\operatorname{coker}(L_i\to M_1)\to
M_2\to L_{i+4}\to0$ に $\Hom_D(-,L_b)$ を適用すると、
$\Hom_D(M_2,L_b)$ から $\Hom_D(M_1,L_b)$ への写像は単射である。
像は表の kernel に含まれ、その次元は表の数と一致する。
従って窓の各頂点で中間は完全である。
頂点 $i$ では対角の $k$ が augmentation として残り、
すべての微分は異なる頂点間の成分からなるので解消は最小である。
三回のmutationは頂点1,2,3を順次消す
Yoneda truncationなので、
$D$上に実現すると最後の$L_{i+4}[-2]$のみが残る。
\end{proof}

\begin{proposition}[spherical helixによるAS正則性]
\label{prop:qa-as-helix-reconstruction}
任意の曲線型rotated square avatarから
$B=\End_{\rm dir}(L_0,L_1,L_2,L_3)$を作る。
その標準射影対象を$E_0,\ldots,E_3$とし、
$E_{i+4}=S_B^{-1}(E_i)[2]$で延長すると
\[
 A(D,P)_{ji}=\Hom(E_i,E_j)\quad(j\ge i)
\]
は$\mathrm{AS}_{\mathrm{rsq}}$に属する。
\end{proposition}
\begin{proof}
$\omega_D=\OO_D$かつ
$H^0(\OO_D)=H^1(\OO_D)=k$だから
各$L_i$は一次元spherical objectである。
全整数のcanonical完全列と
命題~\ref{prop:qa-core-forward-ext-vanishing}、
前補題はacyclic spherical helixを与える。
このacyclicityは同文献Definition~4の全$i<j$に対する
次数零以外のHom消滅である。
これを入力として、対応するexceptional helixの
前向きHomが次数零に集中することに
\cite[Theorem 5]{OkawaUedaSpherical2020}を用いる。
指定した最小分解の頂点・重複度と双対の同定には、
以下のcanonical複体とSerre traceの計算を使う。

$C$を四線束のfull dg Endとする。
$B$は前向き次数零Homと対角の恒等射を含み、
$B\to C$ は dg の実現写像である。
$H^0(C)$ の後ろ向き成分は前向き Ext 消滅の双対により零で、
$H^0(B)\to H^0(C)$ は同型となる。
その cone の唯一のコホモロジーは $H^1(C)$ である。
Serre trace による pairing は合成に関して自然であるから、
$H^1(C)=DB$ は $B$ 双加群としての同型である。
単一のコホモロジーをもつ cone は双加群の導来圏で
$DB[-1]$ に同型なので、三角形
\begin{equation}
 DB[-2]\longrightarrow B\longrightarrow C
 \longrightarrow DB[-1]
 \label{eq:qa-as-dg-triangle}
\end{equation}
がある。
有限代数$B$のSerre関手を用いると
realizationとその右随伴の間に
$S_B[-2]\to\mathrm{id}\to\iota^*\iota_*$を得る。
前補題とこの三角形から
$E_i$のrealizationは全整数で$L_i$となり、
\begin{equation}
\RHom(E_i,E_{j-4})\longrightarrow
\RHom(E_i,E_j)\longrightarrow
\RHom_D(L_i,L_j)
\label{eq:qa-as-rhom-comparison}
\end{equation}
が成立する。
$0<j-i<4$ では $j-4<i$ であり、
$i$ と $j-4$ は同じ窓に入るので最初の項は零である。
従って実現は各窓の前向き Hom 上で合成を保つ同型を与え、
前補題の窓内の計算を例外対象の窓に移せる。
$j-i=4$ の最初の項は対角の $k$ である。
$j-i>4$ では最初の項に距離についての帰納法を使うと、
最後の項も前向き Ext 消滅により次数零に集中するため、
中央の前向き Hom も同じ性質を持つ。
この比較の向きは \cite[式 (9), (30)--(32)]{OkawaUedaSpherical2020}
の $E_{j-4}=S_B(E_j)[-2]$ と一致する。

四項複体を得るため、まず窓 $[i-4,i-1]$ を取る。
前補題の共変 Yoneda 解消の三つの射影項をこの窓に戻すと、
$E_{i-4}$ から始まる逐次右 mutation の複体になる。
実際、三項の 畳み込み は後三対象への導来 Hom が零であり、
$E_{i-4}$ への射の cone は後三対象で生成される。
従って窓の半直交分解によって一意に定まる右 mutation である。
窓の Serre 関手による同定から、その最後の項を $E_i$ とした
四項複体 $K_i$ の全複体は零である。
$K_i$ は次数 $-3,-2,-1,0$ に置き、両端は
$E_{i-4},E_i$ とする。中間項の重複度は前補題の解消から
定義~
ef{def:qa-as-square-rotated-minimal} の rotated 型の表となる。

$RHom(E_j,K_i)$ の項別スペクトル系列で次数零の行を取ると、
右射影複体の頂点 $j$ における評価になる。
$j<i$ では前向き消滅と幅四の窓の半直交性により
他の行は零なので、この行は完全である。
$j=i$ では唯一の例外が左端の
$Ext^2(E_i,E_{i-4})=k$ である。実際、
$S_B(E_i)=E_{i-4}[2]$ なのでこれは Serre 双対から従う。
全複体の消滅により、この $k$ から次数零の行の右端の
コホモロジーへの $d_3$ が同型である。
$j>i$ での評価は directed の定義により零である。
従って右射影複体の唯一のコホモロジーは次数零の $S_i$ である。

双対複体を左頂点 $j$ で評価するには $RHom(K_i,E_j)$ を用いる。
$j>i-4$ では同じ前向き消滅と半直交性により次数零の行が完全であり、
$j<i-4$ では directed の定義により評価は零である。
$j=i-4$ では、次数 $0,1,2,3$ のうち最初の項から来る
$Ext^2(E_i,E_{i-4})=k$ だけが次数零の行の外に残る。
全複体が零なので $d_3$ はこれを次数零の行の
次数三のコホモロジーに同定する。
従って $Ext^3_A(S_i,A)={}_ell S_{i-4}$ であり、
他の次数の Ext は零である。
異なる頂点間の成分しか現れないので右解消は最小である。
\end{proof}


\begin{proposition}[有限Endと全代数の相互逆性]
\label{prop:qa-as-inverse-foundations}
$A\in\mathrm{AS}_{\rm rsq}$から取り出した
$(Y,D,P)$の初期四頂点Endは
\[
A|_{\{0,1,2,3\}}\simeq
\End_{\rm dir}
(P,\OO_D(h-2f),\OO_D(h-f),\OO_D(2h-f))
\]
であり、古典型は$Y$上の可換基準collectionである。
逆にavatarの有限Endから構成した代数の
点表現と余核線束は元の$(Y,D,P)$を回収する。
\end{proposition}
\begin{proof}
後者は命題~\ref{prop:qa-core-directed-foundations}の
二行二列評価行列の局所標準形
$\operatorname{diag}(1,q)$から
$\PP_Y(\coker\kappa)=D$を回収し、
忘れられた線束の正規化で$P$を得ることによる。
前者は、AS代数から作った
\[
0\to\OO_Y^{\oplus2}\xrightarrow{\kappa}
\OO_Y(h)^{\oplus2}\to\iota_*Q\to0
\]
を $\OO_Y(-h)$ で捩ると、$\OO_Y(-h)$ の全コホモロジー消滅により
$H^0(QH^{-1})$ は $0\to2$ の二矢空間と同型になる。
元の列の大域切断は
\[
 0\longrightarrow R_0\longrightarrow
 H^0(Y,\OO_Y(h))\otimes H^0(D,QH^{-1})
 \longrightarrow H^0(D,Q)\longrightarrow0
\]
である。最初の射は普遍評価行列の係数なので、
元の二関係の包含そのものである。
最後の射は可逆層の切断の乗法であり、
$0\to3$ の合成を同定する。

保持した $1,2,3$ の三頂点では、$f,h,h+f$ から $2h$ を引いた
$f-2h,-h,f-h$ のすべてが acyclic なので、
それぞれの切断空間は $Y$ から $D$ への制限で変わらない。
$Y$ 上の $f$ と $h$ の乗法は、
$\PP^1$ 上で $\mathcal E$ の評価列を $\OO(1)$ で捩って得る
全射 $V\otimes H^0(\OO(1))\to H^0(\mathcal E(1))$ であり、
核は $R_1$ である。
従って $1\to3$ の合成も元の二関係で割った積と一致する。
残る
$0\to1$のHomの消滅は
$H^\bullet(\OO_D(h-2f)P^{-1})=0$である。
古典型では $P=\OO_Y$ なので $Q=\OO_Y(2h-f)$ である。
一定階数一の $\kappa$ の像を $I$ とすると、
$0\to I\to\OO_Y(h)^2\to Q\to0$ の行列式から
$I=\OO_Y(f)$、さらに $\ker\kappa=\OO_Y(-f)$ となる。
$\OO_Y^2\twoheadrightarrow I$ の二切断は、
$\OO_Y(f)$ の全生成 pencil の基底をなす。
また $I\hookrightarrow\OO_Y(h)^2$ は、
$\OO_Y(h-f)$ の全生成 pencil の基底を双対化して得る包含である。
いずれも二切断が線形従属なら、正次数の線束を
一切断で生成することになるため不可能である。
この分解で $H^0(\OO_Y^2)\to H^0(I)$ は同型であり、
$H^1(I)=0$ なので、$0\to3$ の積の核は元の $R_0$ となる。
$1\to3$ の積は曲面上の $\mathcal E$ の評価列から同様に $R_1$ を回収する。
従って古典型でも全合成を含めて可換基準 collection を得る。
全$\mathbb Z$-代数は
補題~\ref{lem:qa-as-serre-unique-recovery}が決めるから
有限Endの同型で止まらない。
\end{proof}

\begin{theorem}[squareとrotated squareのavatar分類]
\label{thm:qa-as-avatar-bijection}
定義~\ref{def:qa-as-square-rotated-minimal}の
最小分解型を持つAS正則$\mathbb Z$-代数について
\begin{equation}
\mathrm{AS}_{\rm rsq}/\cong\simeq
 \mathrm{Av}_{\rm rsq}/\cong,\qquad
\mathrm{AS}_{\rm sq}/\cong\simeq
 \mathrm{Av}_{\rm sq}/\cong .
\label{eq:qa-as-avatar-two-bijections}
\end{equation}
$\mathrm{Av}_{\rm rsq}$は
$D\in|2h|$（$\Sigma_2$では$D\cap S=\varnothing$）と
$P\in\Pic^0(D)\setminus\{\OO_D,\OO_D(h-2f)\}$、
$\mathrm{Av}_{\rm sq}$は順序付きrulingを持つ
$D\subset\PP^1\times\PP^1$の$(2,2)$反標準因子と
$P\in\Pic^0(D)\setminus\{\OO_D\}$からなる。
各々に一つの古典型を含める。
曲線の滑らかさ・被約性・既約性は課さない。
逆構成の有限foundationは
\[
B_{\rm rsq}=\End_{\rm dir}(P,HF^{-2},HF^{-1},H^2F^{-1}),
\qquad
B_{\rm sq}=\End_{\rm dir}(P,F,G,FG).
\]
\end{theorem}
\begin{proof}
rotated型では命題~\ref{prop:qa-as-extraction}が必要性、
命題~\ref{prop:qa-as-helix-reconstruction}が
任意avatarのAS実現、
命題~\ref{prop:qa-as-inverse-foundations}と
補題~\ref{lem:qa-as-serre-unique-recovery}が
両構成の相互逆性を与える。
代数同型は点表現・普遍線束を保ち、
avatar同型はHomと合成を保つので、
頂点を保つ代数同型類の全単射となる。

square型代数を一つ進めた
$A^{[1]}_{ji}=A_{j+1,i+1}$は
表からrotated型であり逆shiftも存在する。
一方定理~\ref{thm:qa-core-avatar-square-rotated}は
二つのavatarの同型類を相互変換する。
rotatedの一つ前の四項を曲線上の共通線束
$FP^{-1}$で正規化すると
$(L_{-1},L_0,L_1,L_2)\otimes FP^{-1}
=(P,F,G,FG)$となる。
これは曲線上の操作であって、
非可換圏上にtensor自己同値を仮定しない。
全代数のSerre復元もshiftと可換なので
square側の全単射が従う。
とくに$P=FG^{-1}$や$P=GF^{-1}$を
squareの許容条件から排除してはならない。
この共鳴はrotatedの$\Sigma_2$への移行を
記録し得るがAS正則性を破らない。
\end{proof}

\begin{proposition}[左右Noether性]
\label{prop:qa-as-avatar-noetherian}
上記二型の全代数は左右Noetherianである。
曲線型には四次の正規正則family
$g_i\in A_{i+4,i}$があり、その商は
$D$のtwisted section algebraを与える。
\end{proposition}
\begin{proof}
式~\eqref{eq:qa-as-rhom-comparison}の次数零から
$j>i$について
\[
0\to A_{j-4,i}\xrightarrow{\,g\,}A_{ji}
\to H^0(D,L_jL_i^{-1})\to0
\]
を得る。$t=S_B^{-1}[2]$ とし、三角形の最初の自然変換を
$\eta:t^{-1}\to\operatorname{id}$ と書く。
$E_i\to E_j$ の射 $a$ に対して自然性は
\[
 a\,\eta_{E_i}=\eta_{E_j}\,t^{-1}(a)
\]
を与える。$t^{-1}$ は各 Hom で同型なので、
$\eta$ を前から掛けた像と後ろから掛けた像は一致する。
上の短完全列から後合成による乗法は単射であり、
この式を使えば前合成による乗法も単射となる。
また自然性を $\eta_{E_j}$ 自身に適用すると
$\eta_{E_j}t^{-1}(\eta_{E_j})=
\eta_{E_j}\eta_{E_{j-4}}$ となる。
後合成の単射性により $t^{-1}(\eta_{E_j})=\eta_{E_{j-4}}$
である。従って $g_i=\eta_{E_{i+4}}$ は四周期の
折り畳みと整合する両側正規正則の族となる。
偏極の移動公式に加えて、正規化した初期窓では
$P=(L_0L_1^{-1})HF^{-2}$ である。
この式を次窓に適用すると
$\sigma^*P=(L_4L_5^{-1})(\sigma^*H)(\sigma^*F)^{-2}=P$
を得る。したがって周期式~\eqref{eq:qa-core-rsq-periodic-L}に
$\sigma^*F=FP^2$、$\sigma^*H=HP^4$、$\sigma^*P=P$ を代入すれば、
全整数で $L_{i+4}=\sigma^*L_i\otimes\OO_D(2h)$ となる。
$2h$ は $\Sigma_2$ でも $D\cap S=\varnothing$ により
各成分上正次数なので $D$ 上 ample である。
曲線上の数値類は各既約成分の次数で測られ、
$\sigma$ はこれらを有限置換する。
従って数値類への作用は準冪単であり、
\cite{Keeler2003,Keeler2018Corrigendum} の判定で
この偏極は $\sigma$-ample かつ $\sigma^{-1}$-ample となる。
twisted section algebra は左右 Noetherian である。
四周期をまとめた代数を $J$ とし、$g$ を四位相における
$g_i$ の和とする。$J/gJ$ の対角四隅は
$\OO_D(2h)$ のねじれ切断環である。
異なる位相の隅は、それぞれの対角隅上で
可逆層の切断加群となる。両方向の ampleness により
これらは左右で有限生成なので、$J/gJ$ は対角環の直積上で
左右有限加群となり、左右 Noetherian である。

正規正則性により、$g$-adic 随伴次数環は
$J/gJ$ 上の skew polynomial 環となる。
実際、各 $g^mJ/g^{m+1}J$ は $J/gJ$ と同定でき、
乗法は $g$ を移す正規自己同型でねじれる。
従ってこの随伴次数環は左右 Noetherian である。
内部次数で斉次な右 ideal について、その $g$-adic 初項の
ideal の有限斉次生成系を持ち上げる。
元の斉次元から初項を消すたびに $g$-adic 位数は増え、
$g$ の内部次数が正なので、その元の次数を固定すれば
有限回で停止する。よってこれらの持ち上げは斉次 ideal を生成する。
一般の右 ideal には内部次数の最高次項を使い、
次数を下げる帰納法で有限生成性を得る。左 ideal も同様である。
最後に $A$ の各射影行・列の斉次部分加群は、
折り畳み後に位相の冪等元を保つ斉次部分加群として扱える。
斉次生成元を元の添字へ戻せるので、$A$ も左右 Noetherian である。
古典型は通常の可換曲面のNoether性に帰着する。
\end{proof}

\begin{proposition}[Noether性からcoherentな導来同値へ]
\label{prop:qa-as-avatar-coherent-morita}
前定理のsquare／rotated square型の代数$A$について
$X=\Qgr A$、$\coh X=\qgr A$と置き、
$T=E_0\oplus E_1\oplus E_2\oplus E_3$、
$B=\End(T)$とする。
すると
\[
 \RHom_X(T,-):
 D^b(\coh X)\xrightarrow{\ \sim\ }D^b(\mathrm{mod}\,B)
\]
は三角圏同値である。
さらに$\coh X$の完全自己同値$\nu$があり、
$\nu E_i\simeq E_{i-4}$、
$D^b(\coh X)$のSerre関手は
$\nu[2]$と同型になる。
これは\emph{導来}Morita同値であり、
$\coh X\simeq\mathrm{mod}\,B$という
Abel圏同値を主張しない。
\end{proposition}
\begin{proof}
命題~\ref{prop:qa-as-avatar-noetherian}より
$A$は左右Noetherianである。
有限生成graded加群の有限表示のsyzygyに、
左右の最小単純分解と$\operatorname{Tor}$の検出を繰り返すと、
射影分解は長さ三以内で停止する。
実際、さらに非零のsyzygyがあれば、
左単純加群との$\operatorname{Tor}_{>3}$に非零の項を持つが、
左単純の長さ三の射影分解に矛盾する。
従って各連接対象は有限個の$E_i$のconeから生成される。
補題~\ref{lem:qa-as-serre-unique-recovery}の
四周期完全列により、これらの$E_i$は
$T$が生成する厚部分圏に入る。
$T$はstrong exceptional tiltingであり、
そのEndは有限大域次元二の$B$である。
これで導来Morita同値が$\coh X$の
有界導来圏全体へ延長される。

四周期のSerre同型
$E_{i+4}\simeq S_B^{-1}(E_i)[2]$を
全整数で自然化した$\mathbb Z$-代数のtransportは、
graded加群のexactな四頂点ずらしと可換で、
torsionを保つ。
従って商圏の完全自己同値$\nu$を与え、
$\nu E_i=E_{i-4}$となる。
有限tilting generatorとその全Endにおいて
$\nu[2]$は$B$のSerre関手と一致する。
生成性からSerre関手として同型となる。
\end{proof}

\begin{remark}[一般定理への依存と未証明の比較]
\label{rem:qa-as-classification-scope}
本節では元PDFの体上分類の
必要性・逆構成・全代数の相互逆性を移管した。
spherical helix実現の一般定理
\cite[Theorem 5]{OkawaUedaSpherical2020}と
twisted section algebraのNoether性は外部入力である。
これはAS代数の同型類の対応であり、
別に定義された内在的O-marked$\HHstack$との
pointed同値や任意の非被約基底上のAS相対分類を
意味しない。
\end{remark}
